\section{Problem Setup}
\label{sec:preliminaries}
We consider a target architecture with one software processor and a hardware substrate with limited total area. Software tasks execute sequentially on the CPU, while hardware tasks can overlap subject to precedence constraints. Communication between tasks assigned to different platforms incurs an edge-dependent delay. The goal of HW/SW partitioning is therefore to determine both a placement and an implied execution order that minimize end-to-end latency under the hardware-area budget.

\subsection{Discrete Problem Formulation}
Let $\gG(\gV,\gE)$ be a directed acyclic task graph (DAG). Each node $i \in \gV$ has software execution cost $t_i^{sw}$, hardware execution cost $t_i^{hw}$, and hardware area cost $A_i$. Each edge $(j,i)\in\gE$ encodes both a precedence relation and a communication cost $C_{ji}$ that is paid when tasks $j$ and $i$ are assigned to different platforms.

We represent a discrete partition by a binary vector $\vp \in \{0,1\}^{|\gV|}$, where $p_i=1$ means that task $i$ is mapped to hardware and $p_i=0$ means that task $i$ remains in software. Since software tasks share one processor, the final makespan also depends on the software order $\pi^{sw}$. Given $(\vp,\pi^{sw})$, let $(\vs,\vf)=\Sched(\vp,\pi^{sw})$ denote the induced start-time and finish-time vectors. The true objective is
\begin{align}
\min_{\vp \in \{0,1\}^{|\gV|},\;\pi^{sw}} \quad
& M(\vp,\pi^{sw}) = \max_{i\in\gV} f_i
\label{eq:true_obj}\\
\text{s.t.}\quad
& \sum_{i\in\gV} p_i A_i \le A_{\max}, \label{eq:true_area}\\
& f_i = s_i + p_i t_i^{hw} + (1-p_i)t_i^{sw}, \qquad \forall i\in\gV, \label{eq:true_finish}\\
& s_i \ge f_j + \mathbb{I}[p_j \neq p_i]\, C_{ji}, \qquad \forall (j,i)\in\gE, \label{eq:true_pred}\\
& s_j \ge f_i, \qquad \forall i,j\in\gV:\; p_i=p_j=0 \text{ and } i \prec_{\pi^{sw}} j, \label{eq:true_sw_serial}\\
& s_i \ge 0, \qquad \forall i\in\gV. \label{eq:true_nonneg}
\end{align}
Equation~\eqref{eq:true_pred} enforces precedence and cross-platform communication delay, while Equation~\eqref{eq:true_sw_serial} enforces single-processor serialization for software tasks.

\subsection{Differentiable Relaxation}
Directly optimizing Equations~\eqref{eq:true_obj}--\eqref{eq:true_nonneg} is difficult because both placement and ordering are discrete. In \diffgnnorder, the binary placement vector is relaxed to $\tilde{\vp}\in[0,1]^{|\gV|}$, where $\tilde{p}_i$ is the confidence that node $i$ should be implemented in hardware. The software order is relaxed through a continuous software-priority vector $\tilde{\vpi}^{sw}\in\mathbb{R}^{|\gV|}$.

Under this relaxation, the execution term becomes
\begin{equation}
\tilde{t}_i(\tilde{\vp}) = \tilde{p}_i t_i^{hw} + (1-\tilde{p}_i)t_i^{sw},
\label{eq:soft_exec_i}
\end{equation}
and the communication term becomes
\begin{equation}
\tilde{c}_{ji}(\tilde{\vp}) = |\tilde{p}_j-\tilde{p}_i|\, C_{ji}.
\label{eq:soft_comm_ji}
\end{equation}
The discrete scheduler is replaced by a differentiable surrogate
\begin{equation}
(\tilde{\vs},\tilde{\vf})=
\widetilde{\Sched}(\tilde{\vp},\tilde{\vpi}^{sw}),
\label{eq:soft_sched}
\end{equation}
which approximates both DAG readiness and software serialization.

Since the max operator is nondifferentiable, the surrogate uses a smooth maximum
\begin{equation}
\SMax_{\beta}(\{a_k\})=\frac{1}{\beta}\log\sum_k e^{\beta a_k},
\label{eq:smax_beta}
\end{equation}
where $\beta>0$ controls approximation sharpness. Let $\Pred(i)$ be the predecessor set of node $i$. The differentiable DAG-readiness term is
\begin{equation}
s_i^{dag}
=
\begin{cases}
0, & \Pred(i)=\emptyset,\\[4pt]
\SMax_{\beta}\!\Big(\{\, \tilde{f}_j + \tilde{c}_{ji}(\tilde{\vp}) \mid j \in \Pred(i)\,\}\Big), & \text{otherwise}.
\end{cases}
\label{eq:soft_rdag}
\end{equation}
The software-readiness term is
\begin{equation}
s_i^{sw}
=
\SMax_{\beta}\!\Big(\{\, \tilde{f}_j + \phi_{sw}(\tilde{\pi}_j^{sw}, \tilde{\pi}_i^{sw}) \mid j \neq i\,\}\Big),
\label{eq:soft_rsw}
\end{equation}
where $\phi_{sw}(\cdot)$ is a differentiable ordering function derived from the learned software priorities.

The relaxed start and finish times are then
\begin{equation}
\tilde{s}_i=\SMax_{\beta}\big(\{s_i^{dag},\,s_i^{sw}\}\big),
\label{eq:soft_start}
\end{equation}
\begin{equation}
\tilde{f}_i=\tilde{s}_i+\tilde{t}_i(\tilde{\vp}),
\label{eq:soft_finish}
\end{equation}
and the differentiable makespan surrogate is
\begin{equation}
\widetilde{M}(\tilde{\vp},\tilde{\vpi}^{sw})
=
\SMax_{\beta}\big(\{\, \tilde{f}_i \mid i \in \Sink(\gG)\,\}\big).
\label{eq:soft_makespan_symbol}
\end{equation}

The relaxed area penalty is
\begin{equation}
\mathcal{L}_{\mathrm{area}}(\tilde{\vp})
=
\left[
\max\left(
0,\;
\sum_{i\in\gV}\tilde{p}_i A_i - A_{\max}
\right)
\right]^2.
\label{eq:l_area}
\end{equation}
The proposed method minimizes the differentiable makespan surrogate together with the area penalty and additional regularization terms described in Section~\ref{sec:method}.
